\DocumentMetadata{
  pdfversion=2.0,
  pdfstandard=ua-2,
  lang=en-US,
  tagging=on,
  tagging-setup={math/setup=mathml-SE,tabsorder=structure}
}
\documentclass[11pt,letterpaper]{article}
\usepackage[margin=0.68in,includefoot]{geometry}
\usepackage{newcomputermodern}
\usepackage{amsmath,amscd,mathtools}
\providecommand{\square}{\mdlgwhtsquare}
\usepackage[hidelinks]{hyperref}
\hypersetup{
  pdftitle={Analysis PhD Exam, May 2011},
  pdfauthor={University of Florida Department of Mathematics},
  pdfsubject={Graduate mathematics examination},
  pdfdisplaydoctitle=true,
  bookmarksopen=true
}
\setcounter{secnumdepth}{0}
\setlength{\parindent}{0pt}
\setlength{\parskip}{0.28em}
\setlength{\emergencystretch}{3em}
\setlength{\leftmargini}{2.15em}
\setlength{\leftmarginii}{2.65em}
\setlength{\leftmarginiii}{3em}
\renewcommand{\labelenumi}{\textbf{\theenumi.}}
\pagestyle{empty}
\raggedbottom
\allowdisplaybreaks
\newenvironment{examproblems}[1]{%
  \begin{enumerate}%
  \setcounter{enumi}{#1}%
  \setlength{\topsep}{0.35em}%
  \setlength{\partopsep}{0pt}%
  \setlength{\itemsep}{0.48em}%
  \setlength{\parsep}{0.18em}%
}{\end{enumerate}}
\newenvironment{examsubparts}{%
  \begin{enumerate}%
  \setlength{\topsep}{0.18em}%
  \setlength{\partopsep}{0pt}%
  \setlength{\itemsep}{0.16em}%
  \setlength{\parsep}{0.1em}%
}{\end{enumerate}}
\newenvironment{examinstructions}{%
  \par\begingroup\itshape
}{%
  \par\endgroup\vspace{0.18em}
}
\makeatletter
\renewcommand\section{\@startsection{section}{1}{0pt}%
  {-1.25ex plus -.25ex minus -.1ex}%
  {0.55ex plus .1ex}%
  {\normalfont\large\bfseries}}
\renewcommand{\@maketitle}{%
  \newpage\null\vskip -1.2em%
  {\footnotesize\bfseries
    \href{https://www.ufl.edu/}{University of Florida}, \href{https://math.ufl.edu/}{Department of Mathematics}\par}%
  \vskip 0.18em%
  \tagstructbegin{tag=Title}%
    {\LARGE\bfseries\href{https://gma.math.ufl.edu/exams/analysis-phd/}{Analysis PhD Exam}, May 2011\par}%
  \tagstructend%
  \vskip 0.35em%
  \tagmcbegin{artifact}%
    \hrule height 1pt%
  \tagmcend%
  \vskip 0.55em%
}
\makeatother
\title{Analysis PhD Exam, May 2011}
\author{}
\date{}
\begin{document}
\maketitle
\thispagestyle{empty}
\begin{examinstructions}
Answer SIX questions. Write solutions in a neat and logical fashion, giving complete reasons for all steps and stating carefully any substantial theorems used.
\end{examinstructions}
\begin{examproblems}{0}
\item
This problem has two parts.
\begin{examsubparts}
\item[\textnormal{(i)}]
State the Hahn–Banach theorem.
\item[\textnormal{(ii)}]
Show that there exists an isometric linear map from each separable normed space into \(\ell_\infty\).
\end{examsubparts}
\item
Prove that the spaces~\(c\) (of all convergent scalar sequences) and \(c_0\) (of all null scalar sequences) are not isometrically isomorphic when equipped with the sup norm.
\item
State the Closed Graph Theorem and the Banach Isomorphism Theorem; deduce one of these from the other.
\item
Let the sequence \((T_n)_{n=1}^\infty \subset L(X,Y)\) be bounded in operator norm and assume~\(Y\) complete. Prove that 
\[
Z=\{z\in X:(T_nz)_{n=1}^\infty\text{ converges}\}
\]
 is a closed subspace of~\(X\).
\item
Let~\(X\) and~\(Y\) be closed subspaces of a Hilbert space. Prove that if \(X\perp Y\), then \(X+Y\) is closed.
\item
Let \((\Omega,\mathcal F,\mu)\) be a measure space. Define \(\overline{\mathcal F}\) to comprise all those \(A\subseteq\Omega\) for which there exist \(L,U\in\mathcal F\) such that \(L\subseteq A\subseteq U\) and \(\mu(U\setminus L)=0\) and then define \(\overline\mu(A)\) to be the common value \(\mu(L)=\mu(U)\). Show that \((\Omega,\overline{\mathcal F},\overline\mu)\) is a measure space that is complete in the sense that each subset of a null set is null.
\item
Let \((\Omega,\mathcal F,\mu)\) be an arbitrary measure space; let \(p,q,r>1\) satisfy \(\frac1p+\frac1q=\frac1r\). Prove that if \(f\in L_p(\mu)\) and \(g\in L_q(\mu)\) then \(fg\in L_r(\mu)\) with 
\[
\lVert fg\rVert_r\leq \lVert f\rVert_p\lVert g\rVert_q.
\]
\item
Let \(1\leq p<q<\infty\). Show that \(L_p(\mu)\nsubseteq L_q(\mu)\) iff \((\Omega,\mathcal F,\mu)\) contains measurable sets of arbitrarily small positive measure.
\end{examproblems}
\end{document}
